% SGA 3, Expose VI_B, Loop 1 English reconstruction.
% Scope: Section 2 heading and paragraph 2.0 through Corollary 2.4 in full.
% Authority: Polo--Gille Exp6B-13oct24.pdf, component-local pages 8--11
% (combined-reader pages 349--352; original-margin printed pages 329--333).
% OCR and the Reinhold English Markdown are comparison witnesses only.

\section*{2. 局部有限呈示的群及群态射的``开性质''}
\phantomsection\label{sec:sga3-VIB-2}
\addcontentsline{toc}{section}{2. 局部有限呈示的群及群态射的``开性质''}

\medskip
\phantomsection\label{par:sga3-VIB-2.0}%
\noindent\textbf{2.0.---}
以下始终用 \(S\) 表示任意概形; 将 \(S\)-群概形称为 \(S\)-群.
给定 \(S\)-群 \(G\), 用 \(\epsilon\) 表示单位截面, 用
\[
  c:G\longrightarrow G
\]
表示取逆态射, 并用
\[
  \mu:G\times_S G\longrightarrow G
\]
表示乘法态射. 对每个 \(S\)-概形 \(X\), 用 \(\pi\) 或 \(\pi_X\)
表示结构态射
\[
  X\longrightarrow S.
\]

给定 \(S\)-概形态射
\[
  u:X\longrightarrow Y,
\]
的一个性质 \(\mathscr P(u)\).
若只要 \(u\) 满足 \(\mathscr P(u)\), 对每个 \(S\)-态射 \(Y'\to Y\),
态射 \(u_{(Y')}\) 也满足该性质, 则称 \(\mathscr P(u)\)
\emph{在基变换下稳定}.
若 \(\mathscr P(u)\) 满足以下两个条件,
则称 \(\mathscr P(u)\) \emph{关于拓扑 \(\mathcal T\) 为局部性质}
(参见第~IV~篇, \S\S~4 和 6):
\begin{enumerate}
\item[\textup{a)}] \(\mathscr P(u)\) 在基变换下稳定;
\item[\textup{b)}] 只要存在一族 \(S\)-态射 \(Y_i\to Y\),
  它对于拓扑 \(\mathcal T\) 构成覆盖,
  且每个态射 \(u_{(Y_i)}\) 都满足 \(\mathscr P(u)\),
  则 \(u\) 满足 \(\mathscr P(u)\).
\end{enumerate}

\medskip
\phantomsection\label{prop:sga3-VIB-2.1}%
\noindent\textbf{命题 2.1.---}
\emph{设 \(\mathscr P(u)\) 是 \(S\)-概形态射的一个性质,
关于 \(\mathrm{fpqc}\) 拓扑为局部性质.
设 \(u:G\to H\) 为 \(S\)-群态射.
假设 \(G\) 在 \(S\) 上平坦且普遍开.}

\emph{设 \(W\) 为 \(H\) 中使 \(u\) 在其上满足性质
\(\mathscr P(u)\) 的最大开子集, 令 \(V=u^{-1}(W)\). 则
\[
  U=\pi_G(V)
\]
在 \(S\) 中开, 且 \(V\) 是 \(G|_U\) 的开群子概形.}

\(H\) 中使 \(u\) 在其上满足 \(\mathscr P(u)\) 的最大开子集
\(W\) 的存在性, 来自 \(\mathscr P(u)\) 关于 Zariski 拓扑为局部性质.
由于 \(\pi_G\) 普遍开, \(\pi_G(V)\) 是 \(S\) 的开子集 \(U\).
只须证明 \(V\) 是 \(G|_U\) 的群子概形. 因而可假设 \(U=S\).

令
\[
  G'=G\times_SV,\qquad
  H'=H\times_SV,\qquad
  V'=V\times_SV,\qquad
  W'=W\times_SV,
\]
并令 \(u'=u_{(V)}\).
设 \(W'_1\) 为 \(H'\) 中使 \(u'\) 在其上满足
\(\mathscr P(u)\) 的最大开子集.
由于 \(V\) 在 \(S\) 上平坦且普遍开,
\(H'\) 在 \(H\) 上也如此, 下面的
\hyperref[lem:sga3-VIB-2.1.1]{引理~2.1.1} 表明
\[
  W'_1=W'.
\]
现考虑 \(V\)-概形 \(G'\) (分别地, \(H'\)) 的自同构
\(a\) (分别地, \(b\)), 即以对角截面 \(\delta\) 的逆
(分别地, \(u(\delta)\) 的逆) 作右平移, 定义为
\[
  a(g,v)=(g\cdot v^{-1},v),
  \qquad
  \bigl(\text{分别地, }b(h,v)
    =(h\cdot u(v^{-1}),v)\bigr),
\]
其中 \(T\to S\) 是任意态射,
\(g\in G(T)\), \(v\in V(T)\), \(h\in H(T)\) 任取.%
\begingroup
\renewcommand{\thefootnote}{\fnsymbol{footnote}}%
\footnote[1]{\label{note:sga3-VIB-2.1-b-formula}译者注:
法文底本在第二个展示定义中漏掉了 \((h,v)\) 前面的 \(b\).
前一句定义了自同构 \(b\), 紧接着的恒等式也要求所展示的映射为
\(b(h,v)\); 因此这里补上了该符号.}%
\endgroup{}
显然有
\[
  u'\circ a=b\circ u',
\]
这表明 \(W'\) 在 \(b\) 下稳定, 从而 \(V'\) 在 \(a\) 下稳定,
故 \(V\) 是 \(G\) 的群子概形.

\medskip
\phantomsection\label{lem:sga3-VIB-2.1.1}%
\noindent\textbf{引理 2.1.1.---}
\emph{设 \(\mathscr P(u)\) 是 \(S\)-态射 \(u\) 的一个性质,
关于 \(\mathrm{fpqc}\) 拓扑为局部性质.
考虑 \(S\)-概形态射的一个 Cartesian 方形:}
\begin{center}
  \phantomsection\label{diag:sga3-VIB-2.1.1-cartesian-square}%
  \begin{tikzcd}[row sep=large,column sep=large]
    X' \arrow[r,"f'"] \arrow[d] &
    Y' \arrow[d,"g"'] \\
    X \arrow[r,"f"] &
    Y
  \end{tikzcd}
\end{center}
\emph{其中 \(g\) 平坦且开. 设 \(W\) (分别地, \(W'_1\))
为 \(Y\) (分别地, \(Y'\)) 中使 \(f\)
(分别地, \(f'=f_{(Y')}\)) 在其上满足 \(\mathscr P(u)\) 的最大开子集. 则}
\[
  \emph{\(W'_1=W\times_Y Y'\).}
\]

令 \(W'=W\times_Y Y'\). 由于 \(\mathscr P(u)\) 在基变换下稳定,
有
\[
  W'\subset W'_1.
\]
由于 \(g\) 为开映射,
\[
  W_1=g(W'_1)
\]
是 \(Y\) 的开子集. 令
\[
  V_1=f^{-1}(W_1),
  \qquad
  V'_1=V_1\times_{W_1}W'_1;
\]
显然 \(V'_1=f'^{-1}(W'_1)\).
由于 \(g\) 平坦且开, 由 \(g\) 诱导的态射
\[
  W'_1\longrightarrow W_1=g(W'_1)
\]
忠实平坦且开, 因而是 \(\mathrm{fpqc}\) 拓扑的覆盖
(参见 \SGARefLink{sga3:iv:proposition:6:3:1:item:iv:01}{第~IV~篇, 6.3.1(iv)}).
由于 \(f'\) 诱导的态射
\[
  V'_1\longrightarrow W'_1
\]
满足 \(\mathscr P(u)\), \(f\) 诱导的态射 \(V_1\to W_1\)
也满足该性质. 因而
\[
  W_1\subset W,
  \qquad
  W'_1\subset g^{-1}(W_1)
    \subset g^{-1}(W)=W',
\]
故 \(W'=W'_1\).

\medskip
\phantomsection\label{rem:sga3-VIB-2.1.2}%
\noindent\textbf{注 2.1.2.---}
\emph{态射的许多性质关于 \(\mathrm{fpqc}\) 拓扑都是局部性质;
其中包括平坦、(普遍) 开、(局部) 有限型、有限呈示、拟有限
(参见 EGA~IV\(_2\), 2.5.1, 2.6.1 和 2.7.1)、
平滑、平展和非分歧 (EGA~IV\(_4\), 17.7.3) 等性质.
}

\emph{\hyperref[prop:sga3-VIB-2.1]{2.1} 的证明实际上只使用了平坦态射的基变换.
因此, 该命题适用于相对于 \(\mathrm{fpqc}\) 拓扑满足
\hyperref[par:sga3-VIB-2.0]{2.0} 的条件 \textup{b)},
且在平坦态射的基变换下稳定的性质
(例如, 拟紧且支配这一性质).
}

\emph{当然, 对于关于比 Zariski 拓扑更细的拓扑 \(\mathcal T\)
为局部性质的性质, 也可以陈述类似的命题;
此时 \(G\) 所需满足的条件是: \(\pi_G\) 普遍开,
且为拓扑 \(\mathcal T\) 的覆盖.
}

\emph{特别地, 若 \(G\) 在 \(S\) 上平坦且局部有限呈示,
则对在平坦且局部有限呈示态射的基变换下稳定,
并相对于 \(\mathrm{fpqc}\) 拓扑满足
\SGARefLink{par:sga3-VIB-2.0}{2.0} 的条件 \textup{b)} 的性质,
有类似陈述 (例如, 正则、约化、Cohen--Macaulay 等性质;
参见 EGA~IV\(_2\), 6.8).}

\medskip
\phantomsection\label{prop:sga3-VIB-2.2}%
\noindent\textbf{命题 2.2.---}
\emph{设 \(G\) 和 \(H\) 是两个 \(S\)-群,
\(u:G\to H\) 为 \(S\)-群态射. 则:}%
\footnote{\label{note:sga3-VIB-19}编者注: 在
\textup{\SGARefLink{prop:sga3-VIB-2.2:item:i}{(i)}} 的情形中,
开子集 \(V\) 由 \(G\) 中使 \(u\) 平滑、分别地平展、分别地有限呈示且平坦的
所有点组成; 参见编者\hyperref[note:sga3-VIB-6]{注~6}.
而在 \textup{\SGARefLink{prop:sga3-VIB-2.2:item:ii}{(ii)}} 的情形中,
\(V\) 是包含在 \(G\) 中使 \(u\) 普遍开的点集 \(E\) 中的最大开子集,
但 \(E\) 未必开, 如下例所示 (EGA~IV\(_3\), 14.1.3(i)).
设 \(k\) 为域, \(H=S=\operatorname{Spec}A\), 其中 \(A=k[x]\),
并令 \(G\) 为 \(S\)-群 \(\operatorname{Spec}A[y]/(xy)\).
则 \(E\) 是 \(G\) 的单位截面, 它不是开集.}
\begin{enumerate}
\item[\textup{\SGAInlineRefTarget{prop:sga3-VIB-2.2:item:i}(i)}]
  \emph{假设 \(G\) 或 \(H\) 在 \(S\) 上平坦,
  且 \(G\) 或 \(H\) 在 \(S\) 上局部有限呈示.
  令 \(V\) 为 \(G\) 中使 \(u\) 在 \(V\) 上的限制平坦且局部有限呈示
  (分别地, 平滑或平展) 的最大开子集.
  则 \(U=\pi_V(V)\) 在 \(S\) 中开,
  且 \(V\) 是 \(G|_U\) 的开群子概形.}
\item[\textup{\SGAInlineRefTarget{prop:sga3-VIB-2.2:item:ii}(ii)}]
  \emph{假设 \(G\) 或 \(H\) 在 \(S\) 上普遍开,
  令 \(V\) 为 \(G\) 中使 \(u\) 在 \(V\) 上的限制普遍开的最大开子集.
  则 \(U=\pi_V(V)\) 在 \(S\) 中开,
  且 \(V\) 是 \(G|_U\) 的开群子概形.}
\end{enumerate}

先证明 \textup{\SGARefLink{prop:sga3-VIB-2.2:item:i}{(i)}}.
我们说明 \(\pi_G\) 在 \(V\) 上的限制 \(\pi_V\) 平坦且局部有限呈示:
\begin{enumerate}
\item[\textup{a)}] 若 \(\pi_G\) 平坦 (分别地, 局部有限呈示), 则 \(\pi_V\) 也如此;
\item[\textup{b)}] 若 \(\pi_H\) 平坦 (分别地, 局部有限呈示),
  则 \(\pi_V\) 也如此, 因为它是 \(u\) 在 \(V\) 上的限制与 \(\pi_H\) 的复合.
\end{enumerate}
因此, 在陈述所考虑的四种情形中, \(\pi_V\) 平坦且局部有限呈示,
因而普遍开 (EGA~IV\(_2\), 2.4.6).
令 \(U=\pi_V(V)\), 则 \(U\) 在 \(S\) 中开.
只须证明 \(V\) 是 \(G|_U\) 的开群子概形, 可假设 \(U=S\).

令
\[
  G'=G\times_SV,\qquad
  H'=H\times_SV,\qquad
  V'=V\times_SV,
  \qquad\text{并令}\qquad
  u'=u_{(V)}.
\]
由于 \(V\) 在 \(S\) 上平坦且局部有限呈示,
\(H'\) 在 \(H\) 上也如此.
由 EGA~IV\(_4\), 17.7.4,
\(V'\) 是 \(G'\) 中使 \(u'\) 在 \(V'\) 上的限制平坦且局部有限呈示
(分别地, 平滑或平展) 的最大开子集.
按 \SGARefLink{prop:sga3-VIB-2.1}{2.1} 证明中的方式定义自同构 \(a\) 与 \(b\),
则显然 \(V'\) 在 \(a\) 下稳定, 从而 \(V\) 是 \(G\) 的群子概形.

现在证明 \textup{\SGARefLink{prop:sga3-VIB-2.2:item:ii}{(ii)}}.
\(\pi_G\) 在 \(V\) 上的限制 \(\pi_V\) 普遍开:
或因 \(\pi_G\) 普遍开,
或当 \(\pi_H\) 普遍开时, 因其为 \(u\) 在 \(V\) 上的限制与 \(\pi_H\) 的复合.
令 \(U=\pi_V(V)\), 则 \(U\) 在 \(S\) 中开.
只须证明 \(V\) 是 \(G|_U\) 的开群子概形, 可假设 \(U=S\).

与前面相同, 令
\[
  G'=G\times_SV,\qquad
  H'=H\times_SV,\qquad
  V'=V\times_SV,
  \qquad\text{并令}\qquad
  u'=u_{(V)}.
\]
于是 \(\pi_V:V\to S\) 满射且普遍开, 因而 \(G'\to G\) 也如此.
因此, 由 EGA~IV\(_3\), 14.3.4(i),(ii),
\(V'\) 是 \(G'\) 中使 \(u'\) 在 \(V'\) 上的限制普遍开的最大开子集.
与前面相同地定义自同构 \(a\) 和 \(b\),
则显然 \(V'\) 在 \(a\) 下稳定, 从而 \(V\) 是 \(G\) 的群子概形.

\medskip
\phantomsection\label{cor:sga3-VIB-2.3}%
\noindent\textbf{推论 2.3.---}
\emph{设 \(G\) 为 \(S\)-群, \(V\) 为 \(G\) 中在 \(S\) 上
平坦且局部有限呈示 (分别地, 平滑、平展或普遍开) 的最大开子集.
则 \(U=\pi_V(V)\) 在 \(S\) 中开,
且 \(V\) 是 \(G|_U\) 的开群子概形.}

只须在 \(H\) 为单位 \(S\)-群,
\(u\) 为唯一的 \(S\)-群态射 \(G\to H\) 时
应用 \hyperref[prop:sga3-VIB-2.2]{2.2},
因为此时 \(\pi_H\) 是同构, 且
\[
  \pi_G=\pi_H\circ u.
\]

\medskip
\phantomsection\label{cor:sga3-VIB-2.4}%
\noindent\textbf{推论 2.4.---}
\emph{设 \(G\) 为 \(S\)-群.
若单位截面有一个具备以下一种或多种性质的邻域 \(X\):}
\begin{center}
  \emph{\(X\) 在 \(S\) 上平坦且局部有限呈示
  (分别地, 平滑、平展或普遍开),}
\end{center}
\emph{则存在 \(G\) 的开群子概形, 具有相同性质.}

只须应用 \hyperref[cor:sga3-VIB-2.3]{2.3}, 并注意到此处沿用
\hyperref[prop:sga3-VIB-2.2]{2.2} 的记号时, 有
\[
  \epsilon(S)\subset V,
\]
从而 \(U=S\).
